\subsection{应用：III. 连续函数的代数}

引理 4. --- 设 X 是一个紧空间，H 是 Banach 空间 \(\mathcal{C}(\mathrm{X};\mathbf{C})\)（相应地，\(\mathcal{C}(\mathrm{X};\mathbf{R})\)）的一个闭线性子空间。设 a 是 X 的一个具有可数基本邻域系的点；假定对满足 0 < c < d < 1 的任何数 c 和 d 以及 a 的任何开邻域 U，存在一个 \(f \in H\)，使得

\[
| f | \leqslant 1, \quad | f (a) | \geqslant d, \quad | f (x) | \leqslant c \text {   for   all   } x \in X - U.\tag{12}
\]

那么存在一个函数 \(u \in \mathcal{H}\)，使得 \(|u(x)| < |u(a)|\)（对一切 \(x \neq a\)）。

设 \(\left(\mathrm{V}_{n}\right)\) ( \(n \geqslant 1\) ) 是 \(a\) 的一个基本邻域系，并设 \(\lambda, \mu, \varepsilon\) 是满足下列关系的数：

\[
0 <   \lambda <   1, \quad 1 <   \mu <   \mu + \varepsilon \leqslant 1 + \lambda .
\]

于是 \(0 < \lambda / \mu < 1 / \mu < 1\)。我们将对 \(n\) ( \(n \geqslant 1\) ) 用归纳法定义开邻域的一个递减序列 \((\mathrm{U}_n)\)（\(a\) 的），使得 \(\mathrm{U}_n \subset \mathrm{V}_n\)（对一切 \(n\)），以及函数的一个序列 \((h_n)\)（\(\mathcal{H}\) 中的，满足下列关系）：

\[
\left| h _ {n} (x) \right| \leqslant \mu \quad \text { for   all } x \in X\tag{\ensuremath{13_n}}
\]

\[
h _ {n} (a) = 1\tag{\ensuremath{14_n}}
\]

\[
\left| h _ {n} (x) \right| \leqslant \lambda \quad \text { for   all } x \in X - U _ {n}\tag{\ensuremath{15_n}}
\]

\[
\left| \sum_ {j = 1} ^ {n} \lambda^ {j} h _ {j} (y) \right| <   \sum_ {j = 1} ^ {n + 1} \lambda^ {j} \quad \text { for   all } y \in X.\tag{\ensuremath{16_n}}
\]

假定 \(h_m\) 与 \(U_m\)（对 \(1 \leqslant m < n\)）已定义，它们满足前面四个条件（把 \(n\) 换成 \(m\)）；另一方面，令 \(U_0 = X\)。函数 \(\sum_{j=1}^{n-1} \lambda^j h_j\)（当 \(n = 1\) 时等于 0）是连续的，且取值 \(\sum_{j=1}^{n-1} \lambda^j\)（在点 \(a\) 处）；因此存在开邻域 \(U_n\)（\(a\) 的，含于 \(U_{n-1} \cap V_n\)），使得

\[
\left| \sum_ {j = 1} ^ {n - 1} \lambda^ {j} h _ {j} (y) \right| <   \sum_ {j = 1} ^ {n - 1} \lambda^ {j} + \varepsilon \lambda^ {n} \quad \text { for   all } y \in \mathrm{U} _ {n}.\tag{17}
\]

由假设，存在一个函数 \(f \in \mathcal{H}\)，使得

\[
\begin{array}{l} | f (x) | \leqslant 1 \quad \text { for   all } x \in X, \quad | f (a) | \geqslant 1 / \mu , \\ | f (x) | \leqslant \lambda / \mu \quad \text { for } x \in X - U _ {n}. \end{array}
\]

令 \(h_n = f / f(a)\)；于是关系 \((13_n), (14_n)\) 与 \((15_n)\) 成立；令

\[
g = \sum_ {j = 1} ^ {n} \lambda^ {j} h _ {j} = \sum_ {j = 1} ^ {n - 1} \lambda^ {j} h _ {j} + \lambda^ {n} h _ {n}.
\]

由 (17) 与 \((13_{n})\)，对 \(y \in \mathrm{U}_n\) 我们有

\[
| g (y) | <   \sum_ {j = 1} ^ {n - 1} \lambda^ {j} + \varepsilon \lambda^ {n} + \mu \lambda^ {n} \leqslant \sum_ {j = 1} ^ {n + 1} \lambda^ {j},
\]

因为 \(\varepsilon + \mu \leqslant 1 + \lambda\)；对 \(x \in \mathrm{X} - \mathrm{U}_n\)，我们有 \(|h_p(x)| \leqslant \lambda\)（对 \(1 \leqslant p \leqslant n\)），从而也有

\[
| g (x) | \leqslant \sum_ {j = 2} ^ {n + 1} \lambda^ {j} <   \sum_ {j = 1} ^ {n + 1} \lambda^ {j},
\]

这就完成了 \((16_{n})\) 的证明。

这样，级数 \(\sum_{n=1}^{\infty} \lambda^n h_n\) 在 X 中是正规收敛的，因为 \(\lambda < 1\) 且 \(|h_n(x)| \leqslant \mu\)（对一切 \(n\) 及一切 \(x \in \mathbf{X}\)）；用 \(u\) 表示它的和；由于 \(\mathcal{H}\) 是闭的，它属于 \(\mathcal{H}\)。由关系 \((14_n)\)，我们有 \(u(a) = \sum_{n=1}^{\infty} \lambda^n\)；另一方面，若 \(x \neq a\)，则存在一个整数 \(n\) 使得 \(x \notin \mathrm{U}_{n+1}\)；于是 \(|h_{n+k}(x)| \leqslant \lambda\)（对一切 \(k \geqslant 1\)，由关系 \((15_n)\)）；从而，利用 \((16_n)\)，得到

\[
\begin{array}{l} | u (x) | \leqslant \left| \sum_ {j = 1} ^ {n} \lambda^ {j} h _ {j} (x) \right| + \left| \sum_ {j = n + 1} ^ {\infty} \lambda^ {j} h _ {j} (x) \right| <   \sum_ {j = 1} ^ {n + 1} \lambda^ {j} + \lambda \sum_ {j = n + 1} ^ {\infty} \lambda^ {j} \\ = \sum_ {j = 1} ^ {\infty} \lambda^ {j} = | u (a) |. \end{array}
\]

定理 2 (E. Bishop). --- 设 X 是一个紧空间，A 是复 Banach 代数 \(\mathcal{C}(\mathrm{X};\mathbf{C})\) 的一个闭子代数。假定 A 包含常数且分离 X 的点。设 \(a\) 是 X 的一个点；下列条件等价：

\begin{enumerate}
\def\labelenumi{\alph{enumi})}
\item
  存在一个函数 \(f \in \mathcal{A}\)，使得 \(|f(x)| < |f(a)|\)（对一切 \(x \neq a\)）。
\item
  点 \(a\) 是 \(\mathcal{A}\)-极值点且具有可数的基本邻域系。
\end{enumerate}

\(a) \Rightarrow b)\)：设 \(f \in \mathcal{A}\) 满足 \(|f(a)| > |f(x)|\)（对 \(x \neq a\)）；由 No.~4 的命题 9，\(a\) 是 \(\mathcal{A}\)-极值点。另一方面，若 \(\mathrm{U}_n\) 是由 \(x \in \mathrm{X}\)（满足 \(|f(x)| > |f(a)| - 1/n\)）组成之集，则 \(\mathrm{U}_n\) 是 \(a\) 的一个开邻域，且诸 \(\mathrm{U}_n\) 的交化为 \(a\)；由于 \(\mathrm{X}\) 是紧的，诸 \(\mathrm{U}_n\) 构成 \(a\) 的一个基本邻域系（GT, I, §9, No.~1, 定理 1）。

\(b) \Rightarrow a)\)：只须验证 \(b)\) 蕴涵引理 4 的假设。采用该引理的记号，令 \(\varepsilon = \log d / \log c\)；于是 \(0 < \varepsilon < 1\)。由于 \(a\) 是 \(\mathcal{A}_r\)-极值点，存在一个函数 \(g \in \dot{\mathcal{A}}\)，使得

\[
\mathcal {R} (g) \geqslant 0, \quad \mathcal {R} (g (a)) \leqslant \varepsilon , \quad \mathcal {R} (g (x)) \geqslant 1 \text {for} x \in X - U
\]

（No.~3, 命题 6, b)）。令 \(f = c^g\)；由于 \(f\) 是正规收敛级数 \(\sum_{n=0}^{\infty} (\log c)^n g^n / n!\) 的和，我们有 \(f \in \mathcal{A}\)，且

\[
| f | \leqslant 1, \quad | f (a) | \geqslant c ^ {\varepsilon} = d, \quad | f (x) | \leqslant c \text { for } x \in X - U.\tag{Q.E.D.}
\]

推论. --- 再假定 X 是可度量化的。那么下列性质等价：

\begin{enumerate}
\def\labelenumi{\alph{enumi})}
\item
  \(a\) 是 X 的一个 \(\mathcal{A}\)-极值点。
\item
  存在 \(u \in \mathcal{A}\)，使得 \(|u(x)| < |u(a)|\)（对一切 \(x \neq a\)）。
\item
  设 \(\mathfrak{M}\) 是子集 \(M\)（\(X\) 的）所成之集：对每个函数 \(f \in \mathcal{A}\)，\(|f|\) 在 \(X\) 中的上确界至少在 \(M\) 的一点处达到。那么 a 属于所有的集 \(M \in \mathfrak{M}\)。
\item
  设 \(\mathfrak{N}\) 是 X 的这样的子集 N 所成之集：对每个函数 \(f \in \mathcal{A}\)，\(\mathcal{R}(f)\) 在 X 中的上确界至少在 N 的一点处达到。那么 a 属于所有的集 \(N \in \mathfrak{N}\)。
\end{enumerate}

换言之，

\[
\operatorname{Ch} _ {\mathcal {A}} (\mathrm{X}) = \bigcap_ {\mathrm{M} \in \mathfrak {M}} \mathrm{M} = \bigcap_ {\mathrm{N} \in \mathfrak {N}} \mathrm{N}.\tag{18}
\]

由于在可度量化空间中每个点都具有可数的基本邻域系，\(a\) 与 \(b\) 的等价性由定理 2 得出。我们来证明 \(b\) 蕴涵 \(c\)：事实上，\(a\) 是 \(|u|\) 达到其上确界的唯一点；另一方面，\(c\) 蕴涵 \(a\)，因为对每个 \(f \in \mathcal{A}\)，\(\mathrm{Ch}_{\mathcal{A}}(\mathrm{X})\) 与使 \(|f|\) 达到其上确界的点所成之集相交（No.~4, 命题 9）。同样的推理，利用 No.~3 的命题 5，表明 \(d\) 蕴涵 \(a\)）。最后，为看到 \(b\) 蕴涵 \(d\)，我们可以限于 X 不化为单独一点 \(a\) 的情形，因此 \(u(a) \neq 0\)；于是函数 \(v = u / u(a)\) 属于 \(\mathcal{A}\)，且我们有 \(v(a) = 1\) 以及 \(|v(x)| < 1\)（对 \(x \neq a\)），从而 \(\mathcal{R}(v(a)) = 1\) 且 \(\mathcal{R}(v(x)) < 1\)（对 \(x \neq a\)）。由于函数 \(\mathcal{R}(v)\) 仅在点 \(a\) 处达到其上确界，的确有 \(a \in N\)（对每个 \(N \in \mathfrak{N}\)）。

*例. --- 设 \(\mathrm{X}_1\) 是由点 \((z_1, z_2) \in \mathbf{C}^2\)（满足 \(|z_1|^2 + |z_2|^2 \leqslant 1\)）组成之集（\(\mathbf{R}^4\) 中的单位球），并设 \(\mathcal{A}_1'\) 是全纯函数在 \(\mathrm{X}_1\) 上的限制所成之集，这些全纯函数取值于 \(\mathbf{C}\)、定义在 \(\mathrm{X}_1\) 于 \(\mathbf{C}^2\) 中的一个邻域内（所论邻域依所考虑的函数而定）；设 \(\mathcal{A}_1\) 是 \(\mathcal{A}_1'\) 在 \(\mathcal{C}(\mathrm{X}_1; \mathbf{C})\) 中的闭包，它显然是 \(\mathcal{C}(\mathrm{X}_1; \mathbf{C})\) 的一个闭复子代数且分离 \(\mathrm{X}_1\) 的点。对全纯函数应用“极大值原理”表明 \(\mathrm{Ch}_{\mathcal{A}_1}(\mathrm{X}_1)\) 是球面 \(\mathbf{S}_3\)。

在前面的定义中，把 \(X_1\) 换成“多圆柱”\(X_2\)（由关系 \(|z_1| \leqslant 1\) 与 \(|z_2| \leqslant 1\) 定义），这就给出子代数 \(\mathscr{A}_2'\) 与 \(\mathscr{A}_2\)（\(\mathcal{A}_2^{\prime}\) 的闭包，属于 \(\mathcal{C}(\mathrm{X}_2; \mathbf{C})\)）。这里，极大值原理表明 \(\mathrm{Ch}_{\mathcal{A}_2}(\mathrm{X}_2)\) 是由关系 \(|z_1| = 1\) 与 \(|z_2| = 1\) 定义的“环面”。

由这些结果可以推出：不存在 \(X_{1}\) 的一个开邻域到 \(X_{2}\) 的一个开邻域上、且把 \(X_{1}\) 变成 \(X_{2}\) 的解析同构；因为，若 v 是这样一个映射在 \(X_{1}\) 上的限制，则将有 \(A_{2}=vA_{1}v^{-1}\)，于是 v 将把 \(S_{3}\) 变成一个同胚于 \(T^{2}\) 的空间，这是荒谬的，因为 \(S_{3}\) 是单连通的而 \(T^{2}\) 不是。然而人们会注意到，空间 \(X_{1}\) 与 \(X_{2}\) 是同胚的，二者都是 \(R^{4}\) 中具有非空内部的有界凸集。

